\documentclass[conference]{IEEEtran}

\usepackage{graphicx}
\usepackage{booktabs}
\usepackage{amsmath}
\usepackage{hyperref}
\usepackage{multirow}
\usepackage{float}
\usepackage[utf8]{inputenc}

\title{Glaucoma Stage Classification from Retinal Fundus Images Using DenseNet121 Transfer Learning}

\author{
\IEEEauthorblockN{Group 5}
\IEEEauthorblockA{
Deep Learning Course -- USTH B3 \\
University of Science and Technology of Hanoi \\
2026
}
}

\begin{document}

\maketitle

% ============================================================
\begin{abstract}
Glaucoma is a leading cause of irreversible blindness worldwide, and early detection of its progression stage is critical for timely intervention.
This paper presents a DenseNet121-based transfer learning approach for classifying retinal fundus images into three glaucoma severity stages: normal, early, and advanced.
We employ a two-phase training strategy---frozen backbone followed by full fine-tuning---on the publicly available \texttt{moondream/glaucoma-detection} dataset from Hugging Face.
After deduplication to eliminate data leakage across splits, our model achieves \textbf{91.81\%} accuracy, \textbf{0.9065} macro F1-score, and a Quadratic Weighted Kappa (QWK) of \textbf{0.9049} on the held-out test set.
We also investigate the effect of class weighting and enhanced regularization, finding that the baseline configuration without these additions yields superior performance.
These results demonstrate the effectiveness of transfer learning with DenseNet121 for ordinal glaucoma classification.
\end{abstract}

\begin{IEEEkeywords}
Glaucoma classification, transfer learning, DenseNet121, retinal fundus images, deep learning, medical image classification
\end{IEEEkeywords}

% ============================================================
\section{Introduction}
\label{sec:intro}

Glaucoma is a group of eye diseases that damage the optic nerve, often associated with elevated intraocular pressure.
It is the second leading cause of blindness globally, affecting over 80 million people, with projections reaching 111.8 million by 2040~\cite{tham2014global}.
The disease progresses through stages---from normal to early glaucomatous changes to advanced damage---and early detection is essential because vision loss from glaucoma is irreversible.

Retinal fundus imaging is a non-invasive, widely available screening modality.
However, manual interpretation by ophthalmologists is time-consuming and subject to inter-observer variability.
Automated classification systems based on deep learning offer the potential for consistent, scalable screening.

In this work, we train a DenseNet121 model via transfer learning to classify fundus images into three severity stages.
We adopt a two-phase training strategy: first training only the classifier head with a frozen backbone, then fine-tuning the entire network.
We evaluate the model using accuracy, macro F1-score, per-class metrics, and Quadratic Weighted Kappa (QWK)---a metric particularly suited to ordinal classification tasks where the severity of misclassification errors matters.

We also experiment with class weighting and enhanced regularization (increased dropout and weight decay) to address class imbalance and overfitting.
We find that these modifications do not improve performance over the baseline, and we discuss possible reasons.

This report presents a complete evaluation of DenseNet121 for automated glaucoma stage classification, including dataset preprocessing, training methodology, comprehensive results, and ablation studies.

% ============================================================
\section{Related Work}
\label{sec:related}

Deep learning has been widely applied to glaucoma detection from fundus images.
Most prior work frames the problem as binary classification (glaucoma vs.\ normal), achieving high accuracy using various convolutional neural network architectures~\cite{li2018efficacy}.

DenseNet121~\cite{huang2017densely} is particularly well-suited for medical imaging due to its dense connectivity pattern, which encourages feature reuse across layers and enables strong performance with fewer parameters than comparably deep architectures.
Transfer learning from ImageNet-pretrained models has been shown to significantly improve performance on medical imaging tasks where labelled data is limited~\cite{raghu2019transfusion}.

Multi-class staging of glaucoma severity (rather than binary detection) is less commonly studied but clinically more informative, as it directly maps to treatment decisions.
Our work addresses this three-class formulation using the ordinal structure of the disease progression.

Quadratic Weighted Kappa has been used extensively in diabetic retinopathy grading~\cite{kaggle2015dr} and is appropriate here because it penalizes misclassifications proportionally to their distance on the ordinal scale---confusing normal with advanced is penalized more heavily than confusing normal with early.

% ============================================================
\section{Dataset}
\label{sec:dataset}

\subsection{Source and Structure}

We use the \texttt{moondream/glaucoma-detection} dataset hosted on Hugging Face~\cite{hf_glaucoma}.
The dataset contains retinal fundus images labelled with three glaucoma severity classes: \texttt{normal}, \texttt{early}, and \texttt{advanced}.

The dataset is pre-split by the authors into train, validation, and test subsets:

\begin{table}[H]
\centering
\caption{Original dataset split sizes}
\label{tab:original_split}
\begin{tabular}{lc}
\toprule
\textbf{Split} & \textbf{Images} \\
\midrule
Train & 2,847 \\
Validation & 1,259 \\
Test & 1,272 \\
\midrule
\textbf{Total} & \textbf{5,378} \\
\bottomrule
\end{tabular}
\end{table}

\subsection{Data Leakage and Deduplication}

During exploratory analysis, we discovered significant duplication of images across splits: 459 duplicate images shared between train and validation, and 358 between train and test.
Cross-split duplication constitutes data leakage, which would inflate evaluation metrics by allowing the model to memorize test samples during training.

We addressed this by computing MD5 hashes of raw image bytes and removing duplicates within each split.
The deduplicated split sizes are:

\begin{table}[H]
\centering
\caption{Dataset sizes after deduplication}
\label{tab:dedup_split}
\begin{tabular}{lcc}
\toprule
\textbf{Split} & \textbf{Before} & \textbf{After} \\
\midrule
Train & 2,847 & 2,341 \\
Validation & 1,259 & 1,163 \\
Test & 1,272 & 1,221 \\
\bottomrule
\end{tabular}
\end{table}

All experiments in this report use the deduplicated splits.

\subsection{Preprocessing}

All images are processed identically across models:

\begin{enumerate}
    \item \textbf{Resize} to $224 \times 224$ pixels (required by ImageNet-pretrained models).
    \item \textbf{Training augmentation:} random horizontal flip ($p=0.5$), random rotation ($\pm 10°$), and colour jitter (brightness $\pm 0.1$, contrast $\pm 0.1$, saturation $\pm 0.05$).
    \item \textbf{Normalization} using ImageNet channel statistics: mean $= [0.485, 0.456, 0.406]$, std $= [0.229, 0.224, 0.225]$.
    \item \textbf{Label mapping:} \texttt{normal} $\rightarrow 0$, \texttt{early} $\rightarrow 1$, \texttt{advanced} $\rightarrow 2$.
\end{enumerate}

No augmentation is applied during validation or testing.

% ============================================================
\section{Methodology}
\label{sec:method}

\subsection{Model Architecture}

We use DenseNet121~\cite{huang2017densely} pretrained on ImageNet as the feature extractor.
DenseNet's dense connectivity---where each layer receives feature maps from all preceding layers---promotes feature reuse and gradient flow, making it effective for fine-grained visual classification tasks.

The original 1000-class ImageNet classifier is replaced with a single linear layer mapping from 1024 features to 3 output classes:

\begin{equation}
    \hat{y} = W \cdot \text{GAP}(F(x)) + b
\end{equation}

where $F(x)$ denotes the DenseNet121 feature extractor, GAP is global average pooling, and $W \in \mathbb{R}^{3 \times 1024}$.

Total parameters: 6,956,931 (of which 3,075 are in the classifier head).

\subsection{Training Strategy}

We adopt a two-phase transfer learning strategy:

\textbf{Phase 1: Frozen backbone (5 epochs).}
All convolutional layers are frozen; only the classifier head is trained.
This allows the randomly initialized head to converge without corrupting the pretrained features.
A higher learning rate ($10 \times$ base) is used since only the small head is being optimized.

\textbf{Phase 2: Full fine-tuning (3 consecutive runs of up to 20 epochs each).}
All layers are unfrozen and trained end-to-end with the base learning rate.
Each run saves an independent best checkpoint based on validation loss, with early stopping (patience = 5 epochs).

\subsection{Hyperparameters}

\begin{table}[H]
\centering
\caption{Training hyperparameters}
\label{tab:hyperparams}
\begin{tabular}{ll}
\toprule
\textbf{Parameter} & \textbf{Value} \\
\midrule
Optimizer & Adam \\
Base learning rate & $1 \times 10^{-4}$ \\
Phase 1 learning rate & $1 \times 10^{-3}$ \\
Weight decay & $1 \times 10^{-4}$ \\
Batch size & 32 \\
Image size & $224 \times 224$ \\
Loss function & Cross-entropy \\
Early stopping patience & 5 epochs \\
Random seed & 42 \\
\bottomrule
\end{tabular}
\end{table}

\subsection{Evaluation Metrics}

We report the following metrics on the held-out test set:

\begin{itemize}
    \item \textbf{Accuracy:} overall proportion of correct predictions.
    \item \textbf{Macro F1-score:} unweighted mean of per-class F1 scores, treating all classes equally regardless of size.
    \item \textbf{Per-class precision, recall, and F1:} to identify class-specific strengths and weaknesses.
    \item \textbf{Quadratic Weighted Kappa (QWK):} measures agreement between predicted and true ordinal labels, penalizing misclassifications proportionally to the square of their distance on the ordinal scale. Defined as:
\end{itemize}

\begin{equation}
    \kappa_w = 1 - \frac{\sum_{i,j} w_{i,j} \cdot O_{i,j}}{\sum_{i,j} w_{i,j} \cdot E_{i,j}}
\end{equation}

where $w_{i,j} = (i - j)^2 / (N-1)^2$ is the quadratic weight, $O$ is the observed confusion matrix, and $E$ is the expected matrix under random agreement.
QWK ranges from $-1$ (systematic disagreement) to $1$ (perfect agreement), with values above 0.8 indicating strong agreement.

% ============================================================
\section{Results}
\label{sec:results}

\subsection{Training Progression}

\textbf{Phase 1} (frozen backbone, 5 epochs) brought validation accuracy from 64.4\% to 70.4\%, with the best validation loss of 0.7095 at epoch 4.

\textbf{Phase 2} (full fine-tuning) was run in three consecutive blocks.
The best checkpoint was obtained in Run 1, epoch 8, achieving a validation loss of 0.2499 and validation accuracy of 91.6\%.
Subsequent runs (Run 2 and Run 3) showed diminishing returns, with best validation losses of 0.2945 and 0.2599 respectively.

Figure~\ref{fig:training_curves} shows the training curves for Phase 2 Run 1, illustrating the characteristic gap between training and validation loss that widens as the model begins to overfit after epoch 8.

\begin{figure}[H]
\centering
\includegraphics[width=\columnwidth]{results/figures/densenet121_phase2_run1_training_curves.png}
\caption{Training and validation loss curves for Phase 2 Run 1. The model achieves its best validation loss at epoch 8, after which the gap widens.}
\label{fig:training_curves}
\end{figure}

\subsection{Test Set Performance}

Using the best checkpoint from Phase 2, evaluated on the deduplicated test set (1,221 images):

\begin{table}[H]
\centering
\caption{DenseNet121 baseline test set metrics}
\label{tab:test_metrics}
\begin{tabular}{lc}
\toprule
\textbf{Metric} & \textbf{Value} \\
\midrule
Accuracy & 0.9181 \\
Macro F1-score & 0.9065 \\
Quadratic Weighted Kappa & 0.9049 \\
Test loss & 0.2837 \\
\bottomrule
\end{tabular}
\end{table}

\begin{table}[H]
\centering
\caption{Per-class performance on the test set}
\label{tab:per_class}
\begin{tabular}{lcccr}
\toprule
\textbf{Class} & \textbf{Precision} & \textbf{Recall} & \textbf{F1} & \textbf{Support} \\
\midrule
Normal   & 0.9613 & 0.9092 & 0.9345 & 573 \\
Early    & 0.8500 & 0.8500 & 0.8500 & 240 \\
Advanced & 0.9021 & 0.9706 & 0.9351 & 408 \\
\midrule
\textbf{Macro avg} & 0.9045 & 0.9099 & 0.9065 & 1,221 \\
\bottomrule
\end{tabular}
\end{table}

\subsection{Confusion Matrix Analysis}

\begin{figure}[H]
\centering
\includegraphics[width=0.85\columnwidth]{results/figures/densenet121_confusion_matrix.png}
\caption{Confusion matrix on the test set. Diagonal values indicate correct predictions.}
\label{fig:confusion_matrix}
\end{figure}

Key observations from the confusion matrix (Figure~\ref{fig:confusion_matrix}):

\begin{itemize}
    \item \textbf{Advanced glaucoma detection is strongest:} 396 of 408 advanced cases (97.1\%) are correctly identified, with only 4 misclassified as normal. This is clinically important---the model rarely misses severe cases.
    \item \textbf{Normal classification has highest precision:} 96.1\% of images predicted as normal are truly normal, minimizing false reassurance.
    \item \textbf{Early stage is the weakest:} 17 early cases are missed as normal (7.1\%), and 19 are over-classified as advanced (7.9\%). While the 0.85 F1 is the lowest, the nature of errors is partially mitigated: over-classification as advanced leads to further examination rather than missed diagnosis.
    \item \textbf{Very few dangerous misclassifications:} only 4 advanced cases are predicted as normal, representing a 0.98\% rate of the most clinically dangerous error.
\end{itemize}

\subsection{QWK Interpretation}

The QWK of 0.9049 indicates near-perfect ordinal agreement.
This confirms that when the model does make errors, they tend to be between adjacent categories (normal$\leftrightarrow$early, or early$\leftrightarrow$advanced) rather than across the full severity range (normal$\leftrightarrow$advanced).
For clinical deployment, this is a desirable property: the model respects the ordinal structure of glaucoma progression.

% ============================================================
\section{Ablation: Class Weighting and Enhanced Regularization}
\label{sec:ablation}

To address two identified weaknesses---lower early-stage performance and training overfitting---we experimented with:

\begin{enumerate}
    \item \textbf{Inverse-frequency class weights} in the cross-entropy loss, giving higher weight to the underrepresented early class.
    \item \textbf{Increased regularization:} dropout ($p = 0.5$) added to the classifier head, and weight decay increased from $1 \times 10^{-4}$ to $5 \times 10^{-4}$.
\end{enumerate}

\begin{table}[H]
\centering
\caption{Baseline vs.\ regularized configuration comparison}
\label{tab:comparison}
\begin{tabular}{lccc}
\toprule
\textbf{Metric} & \textbf{Baseline} & \textbf{Modified} & \textbf{$\Delta$} \\
\midrule
Accuracy     & 0.9181 & 0.9099 & $-0.0082$ \\
F1 macro     & 0.9065 & 0.8974 & $-0.0091$ \\
QWK          & 0.9049 & 0.8994 & $-0.0055$ \\
\midrule
Normal F1    & 0.9345 & 0.9264 & $-0.0081$ \\
Early F1     & 0.8500 & 0.8327 & $-0.0173$ \\
Advanced F1  & 0.9351 & 0.9330 & $-0.0021$ \\
\bottomrule
\end{tabular}
\end{table}

The modified model performed worse across all metrics.
Contrary to our expectation, class weights \textit{reduced} early-stage F1 from 0.85 to 0.83.

\textbf{Analysis of failure.} Several factors likely contributed:

\begin{itemize}
    \item \textbf{Class weights introduced training instability.} Phase 1 converged more slowly (59.4\% vs.\ 67.7\% final training accuracy), suggesting the weighted loss made the frozen-head optimization landscape more difficult to navigate.
    \item \textbf{The overfitting gap persisted.} Training accuracy still reached 99.5\% while validation plateaued at $\sim$92\%, indicating that the additional regularization was insufficient to constrain the 7-million-parameter model.
    \item \textbf{Single fine-tuning run vs.\ three.} The baseline benefited from three consecutive fine-tuning runs (effectively 34 epochs of fine-tuning with multiple restarts), while the modified model ran a single 20-epoch block.
\end{itemize}

Based on these results, we selected the \textbf{baseline model without class weights or enhanced regularization} as the final DenseNet121 model.

% ============================================================
\section{Discussion}
\label{sec:discussion}

\subsection{Strengths}

The DenseNet121 model achieves strong classification performance with a QWK of 0.90, indicating that its errors respect the ordinal structure of glaucoma progression.
The two-phase training strategy is effective: Phase 1 stabilizes the classifier head before Phase 2 refines the full network.

The model's error profile is clinically appropriate: it has very high recall for advanced cases (97.1\%) and high precision for normal predictions (96.1\%), minimizing the most dangerous failure modes (missing advanced disease or falsely reassuring patients).

\subsection{Limitations}

\begin{itemize}
    \item \textbf{Early-stage detection} remains the weakest point (F1 = 0.85). This is the most clinically valuable stage to detect, as intervention can prevent progression. The early class also has the fewest training samples (smallest support in test set: 240), which likely contributes to lower performance.
    \item \textbf{Overfitting} is evident from the training curves. The model reaches $>$99\% training accuracy while validation accuracy plateaus around 92\%, suggesting the model memorizes training examples rather than learning fully generalizable features.
    \item \textbf{Data quality.} The original dataset contained significant cross-split duplication (459 + 358 images). While we addressed this via deduplication, the underlying data quality raises concerns about label consistency and image provenance.
    \item \textbf{Single dataset, single seed.} Results are from one random seed on one dataset split. Cross-validation or multi-seed evaluation would provide more robust estimates of model performance.
\end{itemize}

\subsection{Future Directions}

Several avenues exist for further improving glaucoma classification performance:

\begin{itemize}
    \item \textbf{Focal loss} to better handle hard examples and class imbalance, particularly for the early stage.
    \item \textbf{Ensemble methods} combining multiple checkpoints or model architectures.
    \item \textbf{Cross-validation} to obtain more robust performance estimates.
    \item \textbf{Explainability analysis} using Grad-CAM or SHAP to visualize which retinal features drive predictions.
\end{itemize}

% ============================================================
\section{Conclusion}
\label{sec:conclusion}

We presented a DenseNet121-based approach for three-class glaucoma stage classification from retinal fundus images, achieving 91.81\% accuracy and a QWK of 0.9049.
The model demonstrates strong ordinal agreement, rarely confusing normal with advanced cases.

Our ablation study showed that class weighting and enhanced regularization did not improve upon the baseline, suggesting that the standard transfer learning pipeline with early stopping is already well-suited to this task and dataset size.

The early-stage class remains the primary area for improvement.
Future work should explore focal loss, more aggressive data augmentation specific to the early class, and XAI methods (Grad-CAM, SHAP) to understand which retinal features drive the model's predictions---particularly for the challenging early-stage cases where the visual differences from normal are most subtle.

% ============================================================
\section*{Reproducibility}

All code is available at \url{https://github.com/USTH-B3-Projects/DL2026-Group5-Project30}.
The dataset is publicly available on Hugging Face (\texttt{moondream/glaucoma-detection}).
Training was performed on an NVIDIA GeForce RTX 5050 Laptop GPU (8GB VRAM) with PyTorch and random seed 42.

% ============================================================
\begin{thebibliography}{9}

\bibitem{tham2014global}
Y.-C.~Tham, X.~Li, T.~Y.~Wong, H.~A.~Quigley, T.~Aung, and C.-Y.~Cheng,
``Global prevalence of glaucoma and projections of glaucoma burden through 2040: A systematic review and meta-analysis,''
\textit{Ophthalmology}, vol.~121, no.~11, pp.~2081--2090, 2014.

\bibitem{huang2017densely}
G.~Huang, Z.~Liu, L.~van~der~Maaten, and K.~Q.~Weinberger,
``Densely connected convolutional networks,''
in \textit{Proc.\ IEEE Conf.\ Computer Vision and Pattern Recognition (CVPR)}, 2017, pp.~4700--4708.

\bibitem{li2018efficacy}
Z.~Li \textit{et al.},
``Efficacy of a deep learning system for detecting glaucomatous optic neuropathy based on color fundus photographs,''
\textit{Ophthalmology}, vol.~125, no.~8, pp.~1199--1206, 2018.

\bibitem{raghu2019transfusion}
M.~Raghu, C.~Zhang, J.~Kleinberg, and S.~Bengio,
``Transfusion: Understanding transfer learning for medical imaging,''
in \textit{Advances in Neural Information Processing Systems (NeurIPS)}, 2019.

\bibitem{kaggle2015dr}
Kaggle,
``Diabetic retinopathy detection,''
2015. [Online]. Available: \url{https://www.kaggle.com/c/diabetic-retinopathy-detection}

\bibitem{hf_glaucoma}
Moondream,
``Glaucoma detection dataset,''
Hugging Face, 2024. [Online]. Available: \url{https://huggingface.co/datasets/moondream/glaucoma-detection}

\end{thebibliography}

\end{document}
